\documentclass[12pt,a4paper]{report}

\usepackage[T1]{fontenc}
\usepackage[utf8]{inputenc}
\usepackage[english]{babel}
\usepackage{geometry}
\usepackage{setspace}
\usepackage{hyperref}
\usepackage{enumitem}
\usepackage{titlesec}
\usepackage{iftex}

\ifPDFTeX
  \usepackage{newtxtext,newtxmath}
\else
  \usepackage{fontspec}
  \setmainfont{Times New Roman}
\fi

\geometry{margin=2.5cm}
\onehalfspacing
\setlist[itemize]{noitemsep, topsep=4pt}

% Requested typography
% Main text: 12pt (documentclass)
% Titles: 14pt
% Big titles (chapters): 16pt
\titleformat{\chapter}[display]
  {\normalfont\bfseries\fontsize{16pt}{20pt}\selectfont}
  {\chaptername\ \thechapter}
  {10pt}
  {}

\titleformat{\section}
  {\normalfont\bfseries\fontsize{14pt}{18pt}\selectfont}
  {\thesection}
  {1em}
  {}

\titleformat{\subsection}
  {\normalfont\bfseries\fontsize{14pt}{18pt}\selectfont}
  {\thesubsection}
  {1em}
  {}

\hypersetup{
  colorlinks=true,
  linkcolor=blue,
  urlcolor=blue,
  pdftitle={PFE Internship Report - Chapters 1 and 2},
  pdfauthor={Your Name}
}

\newcommand{\imageplaceholder}[2]{
\begin{figure}[htbp]
  \centering
  \fbox{\parbox[c][#2][c]{0.86\textwidth}{\centering #1}}
  \caption{#1}
\end{figure}
}

\newcommand{\diagramplaceholder}[2]{
\begin{figure}[htbp]
  \centering
  \fbox{\parbox[c][#2][c]{0.86\textwidth}{\centering Diagram Placeholder: #1}}
  \caption{#1}
\end{figure}
}

\begin{document}

\begin{titlepage}
  \centering

  \begin{minipage}{0.45\textwidth}
    \centering
    \fbox{\parbox[c][2.7cm][c]{0.9\textwidth}{\centering Institution Logo}}
  \end{minipage}
  \hfill
  \begin{minipage}{0.45\textwidth}
    \centering
    \fbox{\parbox[c][2.7cm][c]{0.9\textwidth}{\centering Company Logo}}
  \end{minipage}

  \vspace{2.2cm}
  {\fontsize{16pt}{20pt}\selectfont \textbf{Internship Report (PFE)}\par}
  \vspace{1cm}
  {\fontsize{16pt}{20pt}\selectfont \textbf{Collectra Project}\par}
  \vspace{0.6cm}
  {\fontsize{14pt}{18pt}\selectfont Chapters 1 and 2\par}

  \vfill

  \begin{flushleft}
    \textbf{Student:} Your Full Name \\
    \textbf{Institution:} Your Institution Name \\
    \textbf{Host Organization:} Company / Lab Name \\
    \textbf{Academic Year:} 2025--2026
  \end{flushleft}

  \vfill
  {\large \today\par}
\end{titlepage}

\tableofcontents
\clearpage

\chapter{General Presentation and Project Framing}
\label{chap:general-presentation}

\section*{Introduction}
Today, companies need better digital tools to manage customer debts. In many cases, teams still use Excel files and manual steps. This causes delay, confusion, and missing information.

This chapter explains the context of the Collectra project done during my PFE internship. It presents the company context, the problem, the project goals, the scope, and the main technologies.

\section{Host Company Presentation}
\subsection{Company Identity and Positioning}
I did my internship at \textbf{Nexus}. The company works in \textbf{IT} and focuses on \textbf{AI and Blockchaine}. It wants to improve daily work by using clear digital processes and reliable data.

\imageplaceholder{Place Company Logo Here}{4cm}

\subsection{Brief Company Introduction}
\textbf{Nexus} was created in \textbf{2015}. It has experience in \textbf{IT}. The company gives importance to service quality, fast execution, and continuous improvement.

During the internship, the company provided:
\begin{itemize}
  \item a real work environment;
  \item access to debt follow-up processes;
  \item support and feedback to align the application with business needs.
\end{itemize}

\subsection{Company Organizational Chart (Organigramme)}
The following diagram shows the main departments and reporting lines in the company.

\diagramplaceholder{Company Organizational Chart (Organigramme)}{7cm}

\section{General Context}
Debt follow-up is often difficult when teams use separate tools. Common problems are:
\begin{itemize}
  \item customer and debt data in many files;
  \item no clear history of calls and reminders;
  \item poor coordination between team members;
  \item risk of mixing data between teams or clients.
\end{itemize}

Collectra solves these issues with a SaaS platform based on workspaces, secure access, and standardized actions.

\imageplaceholder{Place Figure 1.1: Existing process or business context diagram}{5cm}

\section{Project Overview}
\subsection{Project Nature}
Collectra is a full-stack web platform used to centralize:
\begin{itemize}
  \item customer records;
  \item debt records;
  \item payment promises;
  \item follow-up actions and history;
  \item team roles and workspace members.
\end{itemize}

The solution uses a modular monorepo architecture with clear separation between frontend, backend API, and database layer.

\subsection{Functional Positioning}
In daily operations, the platform provides:
\begin{itemize}
  \item authenticated dashboard access for internal users;
  \item workspace-level isolation for multi-tenant operation;
  \item team administration through invitation and role assignment;
  \item secure and documented APIs for business features.
\end{itemize}

\imageplaceholder{Place Figure 1.2: High-level platform overview screenshot}{5cm}

\diagramplaceholder{Chapter 1 Process Diagram: Current vs Target Debt Follow-up Workflow}{6cm}

\section{Problem Statement}
The main problem is:

\begin{quote}
\textit{How can we build a secure and scalable multi-tenant platform that helps teams manage debt follow-up, while keeping data separated and all actions traceable?}
\end{quote}

This problem includes technical needs (security, architecture, consistency) and business needs (team collaboration and process control).

\section{Project Objectives}
\subsection{General Objective}
Build a SaaS solution named Collectra to centralize debt management with strong security and clear workspace governance.

\subsection{Specific Objectives}
The specific objectives are:
\begin{itemize}
  \item implement secure authentication and protected APIs;
  \item enforce workspace context and tenant isolation;
  \item support debt lifecycle operations and customer follow-up;
  \item provide team management with MANAGER and AGENT roles;
  \item guarantee data consistency through a relational model and Prisma ORM;
  \item establish a maintainable and extensible technical foundation.
\end{itemize}

\section{Project Scope}
\subsection{In-Scope Features}
This version of the project includes:
\begin{itemize}
  \item user authentication and account access;
  \item workspace creation, selection, and member management;
  \item campaign and debt management workflows;
  \item CSV-based campaign import with row-level validation;
  \item public debt consultation through secure personal links.
\end{itemize}

\subsection{Out-of-Scope at This Stage}
The following points are not included in this stage:
\begin{itemize}
  \item advanced ERP integrations;
  \item full business intelligence modules;
  \item large-scale multi-region deployment strategy.
\end{itemize}

\section{High-Level Technology Orientation}
Collectra relies on a modern TypeScript stack:
\begin{itemize}
  \item \textbf{Frontend}: Next.js, React, TypeScript, Tailwind CSS;
  \item \textbf{Backend}: Hono.js, TypeScript, Zod validation, OpenAPI contracts;
  \item \textbf{Data}: PostgreSQL with Prisma ORM;
  \item \textbf{Security}: Supabase Auth, JWT, HTTP-only cookies, middleware-based checks;
  \item \textbf{Monorepo}: Turborepo and pnpm workspaces.
\end{itemize}

\imageplaceholder{Place Figure 1.3: Global technical architecture}{5.5cm}

\diagramplaceholder{Chapter 1 Architecture Diagram: Collectra Context Diagram}{6cm}

\section{Development Approach}
Development followed an incremental approach:
\begin{itemize}
  \item identify high-priority business requirements;
  \item model entities and apply migration workflow;
  \item implement secure API endpoints;
  \item integrate frontend modules iteratively;
  \item validate behavior through focused technical tests.
\end{itemize}

\section*{Conclusion}
This chapter explained the project context, the company environment, the problem, and the project goals. The next chapter explains the system requirements in a clear and structured way.

\chapter{Requirements Analysis and System Specification}
\label{chap:requirements-analysis}

\section*{Introduction}
This chapter defines what the system must do and under which constraints. It presents stakeholders, actors, requirements, use cases, and business rules.

\section{Requirement Elicitation Approach}
Requirements were collected from four sources:
\begin{itemize}
  \item business expectations of collection teams;
  \item implemented API routes and validation schemas;
  \item documented operational workflows (CSV import, invitation flow, personal-link flow);
  \item current domain model and data constraints in the database schema.
\end{itemize}

This specification matches the current implemented version and can be extended in future chapters.

\section{Stakeholders and Actors}
\subsection{Primary Stakeholders}
\begin{itemize}
  \item \textbf{Business Management}: needs better control, reporting, and performance;
  \item \textbf{Operational Team}: needs reliable tools for daily follow-up;
  \item \textbf{Customers}: need easy and secure access to debt information;
  \item \textbf{Technical Team}: needs maintainable and secure architecture.
\end{itemize}

\subsection{System Actors}
\begin{itemize}
  \item \textbf{Owner/Manager}: manages workspace and team members;
  \item \textbf{Agent}: performs collection tasks and updates debt follow-up;
  \item \textbf{Customer}: views debt data through a personal link;
  \item \textbf{Authentication Provider}: manages user identity and sessions.
\end{itemize}

\diagramplaceholder{Figure 2.1: Use Case Diagram}{5.5cm}


\section{Product Backlog}
The product backlog is a prioritized list of features to implement. Each user story starts with the same sentence pattern and includes a complexity level.

\begin{table}[htbp]
\centering
\renewcommand{\arraystretch}{1.35}
\begin{tabular}{|c|p{9.5cm}|c|}
\hline
\textbf{Priority} & \textbf{User Story} & \textbf{Complexity} \\
\hline
1 & As a manager, I want to authenticate and access the secure dashboard. & Medium \\
\hline
2 & As a manager, I want to create a workspace to isolate my team data. & Medium \\
\hline
3 & As a manager, I want to invite team members and assign roles in the workspace. & High \\
\hline
4 & As an agent, I want to import a CSV file of debts to create a campaign. & High \\
\hline
5 & As an agent, I want to view, create, and update debt records. & Medium \\
\hline
6 & As an agent, I want to record follow-up actions and payment promises. & Medium \\
\hline
7 & As a manager, I want to generate a secure personal link for a customer. & Medium \\
\hline
8 & As a customer, I want to view my debt through a personal link without login. & Low \\
\hline
\end{tabular}
\caption{Collectra Product Backlog}
\label{tab:product-backlog}
\end{table}



\section{Main Use Cases}
\subsection{Import CSV as Campaign}
\textbf{Actor:} Manager or Agent

\textbf{Preconditions:}
\begin{itemize}
  \item Authenticated internal user;
  \item Active workspace selected;
  \item Valid CSV file available.
\end{itemize}

\textbf{Main Flow:}
\begin{enumerate}
  \item User uploads CSV and optional campaign metadata.
  \item System parses and validates file structure.
  \item System validates rows and marks invalid entries.
  \item System persists campaign and valid debts.
  \item System returns detailed import statistics.
\end{enumerate}

\textbf{Postconditions:}
\begin{itemize}
  \item Campaign and valid debts are available in workspace scope.
\end{itemize}

\diagramplaceholder{UC-01 Activity Diagram: CSV Import Workflow}{6cm}

\subsection{Generate Customer Personal Link}
\textbf{Actor:} Manager or Agent

\textbf{Preconditions:}
\begin{itemize}
  \item Authenticated and authorized user;
  \item Target debt belongs to active workspace.
\end{itemize}

\textbf{Main Flow:}
\begin{enumerate}
  \item User requests link generation.
  \item System verifies workspace ownership.
  \item System signs token with expiration.
  \item System returns personal URL and expiration details.
\end{enumerate}

\diagramplaceholder{UC-02 Sequence Diagram: Personal Link Generation}{6cm}

\subsection{Consult Debt via Public Link}
\textbf{Actor:} Customer

\textbf{Preconditions:}
\begin{itemize}
  \item Customer has a valid personal link.
\end{itemize}

\textbf{Main Flow:}
\begin{enumerate}
  \item Customer opens link in browser.
  \item System validates signature and expiration.
  \item System returns public debt information.
  \item Page displays debt details.
\end{enumerate}

\textbf{Alternative Flow:}
\begin{itemize}
  \item Invalid or expired token returns a denied access state.
\end{itemize}

\diagramplaceholder{Figure 2.2: Sequence Diagram for Public Debt Consultation}{5.8cm}

\section{Business Rules and Constraints}
\begin{itemize}
  \item Every private operation must execute in an active workspace context.
  \item Each debt is linked to one campaign and one client.
  \item Campaign, client, and debt must stay in the same tenant scope.
  \item Invitation and role assignment must follow manager-level policy.
  \item Public access depends on token validity and expiration.
  \item Import rows without required data are rejected with reasons.
\end{itemize}

\section{Conceptual Domain Model (Textual)}
The conceptual model of Collectra includes:
\begin{itemize}
  \item \textbf{Workspace} as the primary tenant boundary;
  \item \textbf{User} and \textbf{WorkspaceMember} for role-based access;
  \item \textbf{WorkspaceInvitation} for team onboarding workflow;
  \item \textbf{Campaign} grouping related debt records;
  \item \textbf{Client} as the debtor profile;
  \item \textbf{DebtRecord} as the core financial tracking unit;
  \item \textbf{PaymentPromise} for commitment tracking;
  \item \textbf{CustomerActionHistory} for operational traceability.
\end{itemize}

\diagramplaceholder{Figure 2.3: Conceptual Data Model (ER/UML)}{6cm}

\diagramplaceholder{Figure 2.4: Class Diagram for Core Domain Entities}{6cm}



\section*{Conclusion}
This chapter converted business needs into clear system requirements. It defined actors, required functions, quality constraints, use cases, and business rules. It is the base for the next chapter on detailed design and architecture.

\end{document}
